\documentclass{article}

% arXiv / preprint version (non-anonymized, "Preprint. Work in progress." footer).
% See neurips_guidelines.md for the other modes (review, [final]).
\usepackage[preprint]{neurips_2024}

\usepackage[utf8]{inputenc}
\usepackage[T1]{fontenc}
\usepackage{hyperref}
\usepackage{url}
\usepackage{booktabs}
\usepackage{amsfonts}
\usepackage{amsmath}
\usepackage{nicefrac}
\usepackage{microtype}
\usepackage{xcolor}
\usepackage{graphicx}


\title{Contrastive Forecasting}

\author{%
  Author Name \\
  Affiliation \\
  \texttt{email@example.com}
}


\begin{document}


\maketitle


\begin{abstract}
  TODO: abstract.
\end{abstract}


\section{Introduction}
\label{sec:introduction}

TODO.


\section{Related Work}
\label{sec:related-work}

\paragraph{Time-series forecasting and the foundation-model era.}
Per-series statistical models (ARIMA, exponential smoothing, state-space)
dominated forecasting until global deep models such as
DeepAR~\citep{salinas2020deepar} and N-BEATS~\citep{oreshkin2020nbeats}
showed that one network trained over many series can outperform
per-series fits. Transformers were beaten by simple linear baselines on
these benchmarks until PatchTST~\citep{nie2023patchtst} introduced patch
tokenization: contiguous values are grouped into a single token, the
attention cost drops by the patch length, and the model gains a useful
local prior. TiDE~\citep{das2023tide} pursued an MLP-only line in
parallel, with a residual block over patched inputs that is structurally
close to the encoder we use here.
TimesFM~\citep{das2024timesfm} is the canonical decoder-only patched TS
foundation model: pretrained on roughly 100B time points, its zero-shot
accuracy is on par with supervised specialists on Monash, Darts and ETT.
The follow-ups extend the recipe along different axes:
Chronos~\citep{ansari2024chronos} tokenizes values into a discrete
vocabulary and trains a T5 encoder-decoder over them;
MOIRAI~\citep{woo2024moirai} masks patches and reconstructs them with an
encoder; Time-MoE~\citep{shi2024timemoe} scales to a billion parameters
with sparse routing; Sundial~\citep{liu2025sundial} replaces categorical
heads with continuous flow matching;
Timer-S1~\citep{liu2026timers1} currently leads GIFT-Eval. All of them
forecast in value space, whether the head is regressive, categorical or
generative.

\paragraph{Contrastive representation learning for time series.}
A separate line learns time-series representations by discrimination
rather than reconstruction. Contrastive Predictive
Coding~\citep{oord2018cpc} encodes a series with a strided CNN,
summarises the past with a GRU, and trains the encoder to assign higher
score to the true future latent than to a batch of negatives. The loss,
InfoNCE~\citep{oord2018cpc}, is a tractable lower bound on the mutual
information between the context and the future, so the encoder is
forced to keep whatever in the past predicts the future. The
time-series successors keep the contrastive backbone but rework what
counts as a positive pair: subseries consistency in
T-Loss~\citep{franceschi2019tloss}, temporal neighbourhoods in
TNC~\citep{tonekaboni2021tnc}, contextual consistency under masking and
random crops with hierarchical contrast in
TS2Vec~\citep{yue2022ts2vec}. These methods produce strong
representations for classification and anomaly detection on UCR and UEA,
but no current TS foundation model uses contrastive learning as the
primary pretraining objective.

\paragraph{Latent-space prediction without contrast.}
A third line, originating in vision, predicts in latent space rather
than in input space. The Joint Embedding Predictive
Architecture~\citep{assran2023ijepa} encodes a context block with one
network, encodes the target block with a second network whose weights
are an exponential moving average of the first, and trains a predictor
to map the context latent to the target latent. The loss is L2 in
latent space. Collapse is avoided by the combination of the asymmetric predictor and
the stop-gradient on the EMA target rather than by contrast: there are
no negatives. Two recent papers
carry the idea to time series. LaT-PFN~\citep{verdenius2024latpfn}
combines JEPA with a prior-data fitted network for in-context
forecasting, and TS-JEPA~\citep{ennadir2025tsjepa} adapts the
patch-token JEPA recipe to univariate series, with competitive results
on classification and long-term forecasting. Masked-reconstruction TS
models such as MOMENT~\citep{goswami2024moment} and
MOIRAI~\citep{woo2024moirai} sit closer to JEPA than tokenized or
autoregressive lines do, since they too predict masked content from
visible context, but their target remains the raw values rather than
their representation.

\paragraph{Positioning.}
Decoder-only transformers trained to predict the next token have become
the dominant architecture in language. As Sutskever has argued, this
objective is far less narrow than it looks: predicting the next token
correctly requires the model to encode grammar, style, world knowledge
and arithmetic, since each is a sub-task of the same prediction
problem. We want to recover this property for time series, with a single
decoder-only model that predicts the next patch and learns whatever
structure is useful along the way, but without committing to
value-space prediction. The observed series is a sample from a
distribution generated by a stochastic process, so a value-space loss
penalises the model for failing to predict intrinsic noise that no
model can forecast. Predicting in latent space lets the target be a
representation of the distribution of plausible futures, with the
observed value treated as one of its samples, and the model is not
forced to fit details that carry no semantic content.

Contrastive forecasting is our attempt at this combination. We forecast
the next patch in latent space, as in JEPA, but enforce non-collapse
with explicit InfoNCE~\citep{oord2018cpc} negatives, as in CPC and
TS2Vec. JEPA is a siamese architecture and is therefore at risk of
collapse by construction: the predictor can drive the L2 loss to zero
by mapping every context to the same point. With an InfoNCE loss the
constant function is instead one of the least stable solutions, since
every sample is pushed away from every other sample in latent space,
and the collapsed embedding is actively repelled. Compared to CPC and
its successors, the target is forecasting rather than representation
transfer for classification. To our knowledge, no published work has
measured latent-space forecasting with explicit negatives at TS-FM
scale.


\section{Contrastive Forecasting}
\label{sec:method}

% TODO: opening paragraph (5--7 sentences). State the goal (a
% decoder-only forecaster trained in latent space with an InfoNCE
% objective), summarise the design (encoder + forecaster, two blocks,
% contrastive loss with cross-channel and cross-batch negatives), and
% point the reader to the three subsections below.

\subsection{Architecture}
\label{sec:architecture}

We work with a multivariate time series of shape $T \times C \times F$,
where $T$ is the number of raw timesteps, $C$ is the number of
channels, and $F$ is the number of features per timestep. In the
experiments reported here $F = 1$, but the formalism does not change
when $F > 1$. The series is cut into non-overlapping patches of width
$W$ along the time axis, so that the patched input has shape
$T_p \times C \times W \times F$ with $T_p = T / W$. Each patch is
encoded independently into a latent $h[t, c] \in \mathbb{R}^H$, and we
project it onto the Euclidean unit sphere by $L^2$-normalising, so
that $\|h[t, c]\| = 1$.

The choice of $W$ is dataset-dependent. Values in $\{16, 32, 64, 128\}$
all appear in the recent foundation-model literature. To minimise
compute, we set the stride between consecutive patches to the patch
width, giving non-overlapping windows that cover the full sequence. A
stride of $1$ is also of interest, since it gives one latent per
timestep, provided the loss is adjusted to align each forecast with
the correct future patch.

The model has two blocks. The \emph{encoder} maps a patch to a latent;
the \emph{forecaster} maps the sequence of past latents to a forecast
of the next latent. We write $f[t, c]$ for the forecaster's output at
position $t$ and channel $c$. The training objective is to make
$f[t, c]$ as close as possible to $h[t+1, c]$ under some dissimilarity
$D$,
\[
  D\big(f[t, c],\, h[t+1, c]\big) \to 0.
\]
Because the latents live on the unit sphere we take $D$ to be the
angular distance, and we measure it through its proxy, the inner
product (Section~\ref{sec:loss}).

\subsection{Encoder}
\label{sec:encoder}

The first encoder we tried in this repository is the residual block
used by TiDE~\citep{das2023tide} and reused by
TimesFM~\citep{das2024timesfm}: a two-layer MLP with a ReLU in the
middle, summed with a linear skip from the input, and a layer
normalisation at the output. The input dimension is $W$, the
intermediate dimension is small (we used $64$), and the output
dimension is $H$. This is the canonical input projection of the
patched-decoder line, and we used it as our baseline.

The encoder we currently use is a bidirectional GRU. Each patch is
read as a sequence of $W$ scalar timesteps, which gives the encoder a
chance to model the local temporal structure that an MLP collapses
into a single linear projection. The forward and backward final hidden
states ($128$ units each in our setup) are concatenated and projected
to $H$, summed with a linear skip from the raw patch, and passed
through a layer normalisation.

We have explored only a small fraction of what is possible for the
encoder. Convolutional and shallow self-attention variants are obvious
candidates, and there is no reason in principle to stop there. If
encoding a patch into a useful latent turns out to be a difficult
problem on real-world series, the encoder could itself be a small
transformer. We leave a systematic encoder search as future work.

\subsection{Loss function}
\label{sec:loss}

Let $h, h' \in \mathbb{R}^H$ be two unit-norm latents. We score their
similarity by the inner product $s(h, h') = h \cdot h'$, which is the
cosine of the angle between them. The choice is the cheapest
non-trivial similarity available once the latents are on the unit
sphere: one multiply-add per pair of dimensions, no square root, no
further normalisation at scoring time.

We index each latent and each forecast by the batch element
$b \in [B]$, the patch position $t \in [T_p]$, and the channel
$c \in [C]$. For every anchor $(b, t, c)$, the positive pair is the
forecast at $(b, t, c)$ matched against the future latent at
$(b, t+1, c)$. The negatives we use in production fall into two
categories.
\begin{itemize}
\item \textbf{Cross-channel.} The latent at a different channel of the
  same series at the same time step, $h[b, t, c']$ for $c' \neq c$.
  There are $C - 1$ such candidates per anchor.
\item \textbf{Cross-batch.} The forecast of a different series in the
  batch at the same time step, $f[b', t, c]$ for $b' \neq b$, paired
  against the future latent at $(b, t+1, c)$. There are $B - 1$ such
  candidates per anchor.
\end{itemize}
Cross-channel candidates are computed by a $C \times C$ inner-product
grid per batch element with the diagonal masked out; cross-batch
candidates by a $B \times B$ inner-product grid over the whole batch
with the diagonal masked out, both in single matrix multiplications
that the GPU executes in parallel. In the loss we do not keep the
per-anchor sets separate: at every position $(t, c)$ we aggregate the
candidates of all anchors into two batch-pooled sums, and the same
denominator is used for every anchor at that position. The cross-batch
pool is therefore quadratic in $B$ and dominates the negatives.

For unit-norm $u, v \in \mathbb{R}^H$, write the temperature-scaled
exponential of the inner product as
\[
  \sigma(u, v) \;=\; \exp\!\big(\langle u, v \rangle / \tau\big).
\]
At every position $(t, c)$, the two batch-pooled negative sums are
\begin{align*}
  \mathcal{N}_{\mathrm{xx}}(t, c)
    &\;=\; \sum_{i = 1}^{B} \, \sum_{c' \neq c}
      \sigma\!\big(h[i, t, c],\; h[i, t, c']\big),\\
  \mathcal{N}_{\mathrm{cb}}(t, c)
    &\;=\; \sum_{i \neq j}
      \sigma\!\big(h[i, t+1, c],\; f[j, t, c]\big),
\end{align*}
and the contribution of the anchor $(b, t, c)$ to the loss is
\begin{equation}
  \mathcal{L}_{b, t, c}
  \;=\; -\log \frac{\sigma\!\big(f[b, t, c],\; h[b, t+1, c]\big)}
  {\mathcal{N}_{\mathrm{xx}}(t, c) \;+\; \mathcal{N}_{\mathrm{cb}}(t, c)}.
  \label{eq:loss}
\end{equation}
The denominator depends on $(t, c)$ but not on $b$: every anchor at
the same position shares the same pooled denominator. We average
$\mathcal{L}_{b, t, c}$ over all anchors and use $\tau = 0.07$.

Several other negatives are natural in this setup but are not enabled
in the production loss: past against future, past against forecast,
forecast at $t$ against forecast at $t+1$, and the cross-channel
variants of all three. An earlier version of the loss carried these
terms explicitly; the current production loss does not.

This list is not exhaustive. One can pair latents at $t+1$ with latents
at $t-1$, draw negatives from a queue of past batches to expand the
pool beyond the current batch, or pair channels across different
series in the batch. Different families of time series may benefit
from specific extra negatives tailored to their structure, and we have
not explored this systematically.


\section{Experiments}
\label{sec:experiments}

TODO.


\section{Results}
\label{sec:results}

TODO.


\section{Discussion}
\label{sec:discussion}

TODO.


\section{Conclusion}
\label{sec:conclusion}

TODO.


\bibliographystyle{plainnat}
\bibliography{references}


\appendix

\section{Appendix A}
\label{sec:appendix-a}

TODO.


\end{document}
